\documentclass[11pt]{article}
\usepackage[a4paper,margin=2cm]{geometry}
\usepackage{amsmath,amssymb}
\usepackage[hidelinks]{hyperref}
\allowdisplaybreaks
\title{Qpu.Information.PhysicalQubit.Transmon}
\author{Tsvetan Rouschev}
\date{@uuidna/qpu v1.1.1, CC-BY-NC-ND-4.0}
\begin{document}
\maketitle
\noindent\url{https://github.com/uuidna/qpu/blob/v1.1.1/lean/Qpu/Information/PhysicalQubit/Transmon.lean}

\paragraph{\texttt{Transmon.frequency\_at\_fifty}}
at E\_J = 50 E\_C exactly, f₀₁ = 19 E\_C: √(8 · 50 · E\_C²) = 20 E\_C, for every E\_C
\[ \forall (\mathrm{EC} : \mathbb{N}), (19 \cdot \mathrm{EC} + \mathrm{EC}) \cdot (19 \cdot \mathrm{EC} + \mathrm{EC}) = 8 \cdot (50 \cdot \mathrm{EC}) \cdot \mathrm{EC} \]
\noindent axioms: \texttt{propext}, \texttt{Quot.sound}; proof: \url{https://github.com/uuidna/qpu/blob/v1.1.1/lean/Qpu/Information/PhysicalQubit/Transmon.lean\#L28}

\paragraph{\texttt{Transmon.sycamore}}
Sycamore's measured 6 924 MHz and −208 MHz: 93.1 fF gives E\_C/h = 208 MHz, E\_J/h = 30 569 MHz gives f₀₁, past 50 E\_C
\[ \mathrm{Transmon.charging}\ 93100\ 208 \land \mathrm{Transmon.frequency}\ 30569\ 208\ 6924 \land 50 \cdot 208 < 30569 \]
\noindent axioms: none; proof: \url{https://github.com/uuidna/qpu/blob/v1.1.1/lean/Qpu/Information/PhysicalQubit/Transmon.lean\#L25}
\end{document}
